\section{Pure Literal Elimination}

Suppose I gave you the following CNF equation:

\begin{FVerbatim}
1 2 -3 0
1 -4 5 0
1 -2 5 0
\end{FVerbatim}

Do you notice anything odd about it? Perhaps ``odd" is not the correct word here, do you notice 
an easy way to satisfy this equation? You probably noticed that the variable $1$ shows up in every 
clause and is always in the positive. So, if you just assign $1 = T$, then it doesn't matter what 
the rest of the variables are assigned, because $1 = T$ satisfies every clause! This is idea behind 
\textbf{Pure Literal Elimination}. Essentially, every now and then, we scan our equation for ``pure 
literals'', which are variables that only show up in the positive or negative (but never both), and
then give them whatever assignment is necessary to have their literals evaluate to True. 

Now, suppose I give you this equation:
\begin{FVerbatim}
1 0
-1 -2 -3 0
1 2 0
-2 3 0
\end{FVerbatim}

I have a question for you dear reader: is $-2$ a pure literal? You might be inclined to say ``no"
because $-2$ and $2$ are both present in the equation. However, despite that, the answer is ``yes.''
Why? Well notice in the first clause that $1$ is a unit variable, so $1 = T$ is forced by pure literal 
elimination. Then the third clause, $1 2 0$, because it's already satisfied, the variable $2$
doesn't matter anymore. The other two clauses, which are unsatisfied, are the ones that matter. 
They both contain $-2$, and because they're unsatisfied, we can force $2 = F$ to satisfy them,
thus making the assignment of $3$ irrelevant. 

Okay, that makes sense, sort of. How come we're allowed to do that? Like, how do we know that 
forcing $2 = F$ wouldn't screw up the equation elsewhere (if it were larger)? And the answer is:
because it doesn't hurt the other clauses. In the third clause, once it's satisfied, the only thing 
that can make it unsatisfied is the backtracking, but since $1$ is unit, there shouldn't be any 
backtracking that would suddenly make its satisfiability rest on the assignment of $2$. Likewise, 
for the other two clauses, no matter what value $3$ takes, $2 = F$ must be true to satisfy the 
other clause (try it out for yourself!), so we can go straight ahead and assign $2 = F$ and skip 
deciding. 

Okay, sure, but there are 8 total assignments that we can get. How do we know that forcing 
$2 = F$ would lead to a satisfiable assignment? As in, how do we know we're not prematurely 
boxing ourselves in? That's another great question! The reason why we're not boxing ourselves 
in prematurely is because, again, it only helps the satisfiability of the whole equation, 
and because we're not doing anything detrimental, we can make that decision.
\footnote{Essentially, think of it as making the solution ``more correct'' 
by forcing the pure literal.}

\subsection{Implementing Pure Literal Elimination}

Now that we (sort of) understand why it works, how do we implement it?
We need some form of efficient tracking for things we've already encountered, perhaps 
something like a map or a set. Then, we just need to make sure we've only encountered 
one version of a literal, and if we have, then we know it's pure. Then we just need 
to gather our pure literals and the satisfied clauses, and then return them. 

\begin{algorithm}[H]
\caption{Pure Literal Elimination}\label{alg:15}
\begin{algorithmic}[1]

\Procedure{pureElim}{assignment[1..n]}

\State varToClause \Comment map from variable to clause
\State pure[1..n]
\State satisfiedC[1..m]

\For{$i \gets 1$ up to numClauses}
    \If{satisfied[$i$] = $true$}
        \State move onto the next clause
    \EndIf

    \State $Clause[1..t] \gets Clauses[i]$
    \For{each $var$ in $Clause$}
        \State $\Call{append}{varToClause[var], i}$
    \EndFor
\EndFor

\For{each ($var$, $clauses$) in varToClause}
    \State $idx \gets \Call{abs}{var}$
    \If{assignment[$idx$] $\neq$ Unassigned}
        \State move on to the next pairing
    \ElsIf{$\Call{contains}{varToClause -var}$}
        \State move on to the next pairing
    \EndIf

    \If{$var$ > 0}
        \State $assignment[idx] \gets true$
    \Else
        \State $assignment[idx] \gets false$
    \EndIf

    \For{each $cIdx$ in $clauses$}
        \State satisfied[cIdx] = $true$
        \State $\Call{append}{satisfiedC, cIdx}$
    \EndFor

    $\Call{append}{pure, var}$
\EndFor

\State return (pure, satisfiedC)

\EndProcedure
\end{algorithmic}
\end{algorithm}

So this is a single pass version of pure literal elimination. It is possible that 
after finding some pure literals, and marking some clauses satisfied, a second pass could 
find new pure literals. So why are we implementing a single pass version? Besides 
it not being a technique used in modern SAT solvers, in my experience, which isn't much 
since I'm being honest, it doesn't make a big (positive) difference in solve time 
efficiency. And besides, if we wanted a multi-pass version, we could wrap it in 
\textit{while(true)} loop and break once we find no more new pure literals. 

Now let's get into how the algorithm works. We first build a map from variables 
to their clauses (in unsatisfied clauses). Next, afterwards, we then check 
for each element in the map, that the map does not contain the negative variable 
of the current entry, and if it doesn't, we then need to make sure that the current
variable is Unassigned. If the current entry passes both those conditions (pure and 
Unassigned), we assign it so it evaluates to True. Then, we mark each clause the 
variable is found in as satisfied, and then we add the indices of the satisfied clauses 
into the return values. 

\subsection{Pure Literal Elimination - Results}

So. I have some terrible news. After implementing pure literal elimination, our solving 
time \textit{\textbf{increased}}. I know, I know. It's not supposed to happen, but it can 
happen. There are many possible explanations for this. For instance, pure literal elimination
takes a decent amount of work to perform. We have to scan all unsatisfied clauses, and since 
there is no way to determine if a literal is impure until after everything is scanned
(we could have used a hash set, but then we'd need to hash for every literal), we're 
accruing memory and runtime costs. Second, even if pure literal elimination finds something,
it usually isn't much (at least for the testing sets currently in use), so we're doing 
so much work for so little in return. Third, it increases our backtracking cost because now 
we need to undo unit propagation and the pure literal elimination. Fourth, my implementation
isn't perfect. It is possible that forcing a pure literal may force another pure literal 
in a different clause, so you would need multiple passes to do. However, our implementation 
is just one pass (wrap it in a \textit{while} loop or something).
\footnote{Modern SAT solvers don't tend to use pure literal elimination for reasons outside 
the ``it wasn't very helpful'' aspect, and next chapter we'll see why.}

However, no one cares \textit{why} something performed badly. People only care about 
\textit{how} badly something performed. So here are the results:
For uf20.91: 1548ms, almost a 2x increase from just unit propagation.
For UUF50.218.1000: 231,107ms (~3.85 minutes), about a 1.5 minute increase 
from just unit propagation.

\begin{FVerbatim}
On all 1000 equations in uf20-91:

Pure Literal Elimination:  1548ms (~1.55 seconds).
Unit Propagation:          868ms (~0.87 seconds).
Recursive Backtracking:    1,540,613ms (~25.68 minutes).
Brute Force:               1,770,851ms (~29.51 minutes).
Short-Circuited (Basic):   49,874ms (~50 seconds).
Parallel:                  10,589ms (~10.5 seconds).

On all 1000 equations in UUF50.218.1000:

Pure Literal Elimination: 231,107ms (~3.85 minutes).
Unit Propagation:         140,041ms (~2.34 minutes).
\end{FVerbatim}

Pure Literal Elimination had about a \textbf{2x} increase (which is bad) in solving time
when compared to just Unit Propagation on uf20.91, which is the test set of satisfiable
problems. It also had about a \textbf{1.6x} increase in solving time on UUF50.218.1000,
which is the set of unsatisfiable problems.